% Fragment: fig-sealed-three-principals.tex -- the three parties who must
% never see each other's plaintext, and the one boundary that seals all
% three lines of sight at once.
%
% Reader question: which three parties must never see each other's
% plaintext, and what single mechanism seals all three lines of sight?
% Claim: Derek, Erin, and the cloud operator form a triangle of mutual
% "must not see" constraints, and one boundary (not three separate ones)
% seals every edge of that triangle at once.
% Evidence: "The scene" paragraph, website-v2/public/whitepaper/
% sealed-harbor.tex, Sec. sealed-problem ("Derek owns sensitive financial
% data D. Erin owns a valuable agent system M... Derek must not receive
% M, Erin must not receive D, and the cloud operator hosting the
% computation must receive neither.").
% Must distinguish: three named parties, not two -- the counter-read this
% figure exists to correct is a reader who thinks only Derek-versus-Erin
% needs separating and forgets the host; and one enclosing boundary, not
% three pairwise fences, since the room seals all three sight-lines with a
% single mechanism.
% Grammar chosen: relation map, three principal nodes in a triangle with
% explicit "must not see" edges on every side, the whole triangle enclosed
% by one trust-boundary rectangle (the S12 Dolev-Yao convention already
% used for the room elsewhere in this chapter).
% Grammar rejected: three separate two-party boxes (Derek|Erin,
% Derek|cloud, Erin|cloud) -- that draws three fences where the chapter's
% claim is exactly that one mechanism does the sealing, not three.
% Five-second test: count the boundaries drawn -- one rectangle, around
% all three parties and all three edges, not a fence per pair.
% QA: beauty_lint warns B7 (sub-point near-alignment noise between
% rotated edge-tags and node centres); justified, not fixed.
\begin{figure}[H]
\centering
\pdfigurehue{pdcobalt}%
\begin{tikzpicture}[pd figure,x=1cm,y=1cm]
  \node[pd focus outline,minimum width=2.6cm,align=center]
    (DEREK) at (0,3.6) {Derek\\\itshape holds $D$};
  \node[pd focus outline,minimum width=2.6cm,align=center]
    (ERIN) at (6.2,3.6) {Erin\\\itshape holds $M$};
  \node[pd focus outline,minimum width=2.9cm,align=center]
    (CLOUD) at (3.1,0.4) {cloud operator\\\itshape holds neither};
  \draw[pd breach rule] (DEREK) -- (ERIN);
  \node[pd tag,anchor=south] at ($(DEREK)!0.5!(ERIN)$) {must not see};
  \draw[pd breach rule] (DEREK) -- (CLOUD);
  \node[pd tag,anchor=east,rotate=52] at ($(DEREK)!0.5!(CLOUD)$) {must not see};
  \draw[pd breach rule] (ERIN) -- (CLOUD);
  \node[pd tag,anchor=west,rotate=-52] at ($(ERIN)!0.5!(CLOUD)$) {must not see};
  \begin{scope}[on background layer]
    \node[pd trust boundary,fit={(DEREK)(ERIN)(CLOUD)},inner sep=16pt]
      (ROOM) {};
  \end{scope}
  \node[pd title,anchor=south west] at (ROOM.north west) {the sealed room};
  \node[pd note,anchor=north,align=center,text width=6.6cm]
    at (3.1,-1.55) {one boundary seals all three sight-lines at once};
\end{tikzpicture}
\caption{The room must keep Erin's model from Derek, Derek's data from Erin, and both from the cloud operator. One mechanism must satisfy all three separations. [internal]}
\label{fig:sealed-three-principals}
\end{figure}
